\documentclass[11pt,a4paper]{article}
\usepackage[margin=24mm]{geometry}
\usepackage{amsmath,amssymb,booktabs,array}
\input{glyphtounicode}
\pdfgentounicode=1
\usepackage[hidelinks]{hyperref}
\hypersetup{hypertexnames=false}
\setlength{\parindent}{0pt}
\setlength{\parskip}{5pt}
\newcommand{\pos}[1]{[#1]_+}
\renewcommand{\labelitemi}{\(\bullet\)}
\begin{document}
\begin{center}
{\Large J4: speed advice, actual mismatch and immediate merge work}\\[4pt]
16 October 2026 --- static constitutive qualification, not traffic simulation
\end{center}

\textbf{Status.} J4 is authorized and begun after review of OR-Q1's
conditional recovery evidence. This recheck reaches \textbf{J4-W1}, a
Type 2 G3 interpretation/substitution choice in the trial work law. No
alternative is selected. O-R, R-E, S-E-P, H1/T1, the original action range
and local objective are unchanged. The whole Step 3 remains open.

\section*{1. The law and the physical question}
Retain fixed $w,k_j>0$, $K=wk_j$ and advice $v\in[v_-,v_+]$:
\begin{equation}
C(v)=\frac{Kv}{w+v},\qquad U(g)=\frac{wg}{K-g},\qquad
0<R\leq C(v_-),\quad 0<B\leq R.\tag{1}
\end{equation}
Here $B$ is current shared work, not the far lane. Free marks are $v$;
queued marks are $U(g_j)$. The present trial unit work is
\begin{equation}
c_E^{A}=c_0+\beta\frac{g_A}{C(v)}\frac{\pos{u_A-u_E}}v,
\quad c_0=1+\delta,\quad
\eta=c_0-\beta(1+v_+/w)>0.\tag{2}
\end{equation}
Use $\delta>0,\beta>0$ for strict diagnostics. The previously proved
bound gives a unique queued-E root of
\begin{equation}
H(g,v)=g\left[c_0+\beta\frac{a}{C(v)v}
                      \pos{u_A-U(g)}\right]=x,\quad H_g\geq\eta,
\quad a=g_A,\quad x=B-a.\tag{3}
\end{equation}
When A is free and its offer fits, $a=\lambda_A<B$, $u_A=v$.
When both streams queue, $a=B/2$, $u_A=U(a)$ and $x=a$.

These are local constitutive tests: hold current offers/budget or queues/
budget fixed when comparing advice. They are not two complete H1 trajectories
at different controls, a reachability certificate or a policy experiment.
When both are queued and all relevant flows are below $C(v_-)$, the
free-flow ceilings are inactive; the congested branch $U$ is common.

The physical question is whether a changed but inactive speed sign should
still change immediate interaction work. A direct desired-speed influence
could be posited, but it is an additional constitutive hypothesis. Neither
source-relative-speed evidence nor units alone establishes it.

\newpage
\section*{2. Free A: the speed-matching benefit is not guaranteed}
Put $\rho(v)=1/(C(v)v)$. Then
\begin{equation}
\rho(v)=\frac{w+v}{Kv^2},\quad
\rho'(v)=-\frac{v+2w}{Kv^3},\quad
\frac{d}{dv}\{(v-u)\rho(v)\}
=\frac{(u-w)v+2wu}{Kv^3}.\tag{4}
\end{equation}
In (3), at fixed $g$, $u=U(g)$ is independent of $v$. On the free-A,
queued-E branch $u<v$, so
\begin{equation}
\frac{dg}{dv}=-\frac{H_v}{H_g},\qquad
H_v=\frac{\beta a g}{Kv^3}\{(u-w)v+2wu\}.\tag{5}
\end{equation}
If $u<wv/(v+2w)$, then $H_v<0$ and $dg/dv>0$: lowering advice
\emph{reduces} E service in this local comparison. In other regimes the
sign can be reversed. This is a property of the two changing scales, not
a universal traffic conclusion. Lowering free A speed at fixed flow also
changes its density and physical gaps; their real combined effect is not
proved to be monotone by the user's intuition.

An exact static example, using normalized units, is
\begin{equation}
w=2,\quad k_j=3,\quad v\in[1,2],\quad R=1,\quad
c_0=5/4,\quad\beta=1/10,\quad a=1/2,\quad
B=22183/35400.\tag{6}
\end{equation}
Here $C(1)=2$, $C(2)=3$, $\eta=21/20>0$ and $a<B<R$.
At $v=2$, (3) has the exact root
\begin{equation}
g=1/10,\quad U(g)=2/59,\quad x=B-a=4483/35400.\tag{7}
\end{equation}
At $v=1$, test that same $g$:
\begin{equation}
H(1/10,1)-x=11/14160>0.\tag{8}
\end{equation}
Since $H$ is strictly increasing, its unique new root is smaller than
$1/10$. A is still free with offer $1/2$; the same budget serves it.
Take sufficient existing E backlog for this instantaneous test. No initial
queue, history or parameter set in (6) is adopted as the primary task.

\newpage
\section*{3. Both queued: a direct effect of an inactive advice bound}
Hold $B$ fixed, with both queues positive. S-E-P cooperation gives
$a=B/2$, independent of advice. For $c_0>1$, (3) has
$0<g\leq a/c_0<a$, hence $U(g)<U(a)$ and the positive part is active.
For fixed $g$,
\begin{equation}
H_v=\beta a g\{U(a)-U(g)\}\rho'(v)<0,
\qquad \frac{dg}{dv}>0.\tag{9}
\end{equation}
Changing advice changes E service although A's flow and congested speed,
the fixed FD congested branch and the current budget are unchanged. The
advice bound does not bind either trace. Any later change in E speed or
recovery is a consequence of this extra immediate-work channel.

For example, with $w=2,k_j=3,v\in[1,2],B=R=1,c_0=5/4$ and
$\beta=1/10$, $a=1/2$, $U(a)=2/11<1$. At a fixed E trace, lowering
$v$ from 2 to 1 triples $\rho(v)$ and therefore triples its mismatch
coefficient. Root monotonicity makes E service smaller. This is a static
comparison, not a full-cohort consequence or a proof the state is reachable
with the same initial observations under both controls.

\textbf{Negative controls.} With $\beta=0$ or $a=0$ there is no mismatch
term. With $c_0=1$ and both queued, the unique root is $g=a$, giving equal
traces and zero mismatch, so the strict sign in (9) disappears.

\textbf{Interpretation.} If work is intended to depend only on represented
current physical traffic and an inactive speed cap has no independent
behavioral meaning, (9) fails that intended consistency test. It can instead
be retained as a law with an explicit desired-speed/saturation influence.
That meaning needs evidence. This check is distinct from the technical
invariant-junction condition in the LTM literature; no general node
invariance or hybrid nonexistence theorem is being asserted.

The restricted congested follower calculation in Leclercq et al. (2016),
Section 2, distinguishes insertion speed from free-flow speed. The empirical
relative-speed study does not derive (2). The sign calculations here are
our deductions, not results attributed to those papers.

\newpage
\section*{4. A concrete fixed-scale candidate, not yet selected}
Choose once a physical reference speed $v_\star=v_+$ and
$C_\star=C(v_\star)$, holding both fixed over the original action box.
Replace only (2) by
\begin{equation}
c_E^{F}=c_0+\beta\frac{g_A}{C_\star}
                  \frac{\pos{u_A-u_E}}{v_\star}.
\qquad
\chi=\frac{\beta}{C_\star v_\star}.\tag{10}
\end{equation}
$\chi$ has units s$^2$/(veh-eq m), so $\chi g_A(u_A-u_E)$ is
dimensionless. Reference scales specify a fixed coefficient; changing the
search box must not silently change $\chi$. The anchor choice introduces
no new state and makes (10) identical to (2) at $v=v_+$.

On the regular free-A/queued-E branch with fixed $a,B$ and $U(g)<v$,
\begin{equation}
H_v^{F}=\frac{\beta a g}{C_\star v_\star}>0,
\qquad \frac{dg}{dv}<0.\tag{11}
\end{equation}
Lower actual free A speed then raises E service in this specific
comparison. $U'(g)>0$ raises its speed mark, and the unchanged recovery
profile is nonincreasing in that mark. These are local implications at
fixed other inputs, not a global ordering of recovery histories or total
travel times. H1 arrivals and travel time also change with advice.

With both queued, (3) using (10) has no direct $v$ dependence, so its local
allocation is unchanged while capacities remain nonbinding. The response
depends on actual traffic traces, passing flow and current service instead.

\textbf{Retained sufficient bounds.} Since $C(v)\leq C_\star$ and
$v\leq v_\star$,
\begin{equation}
c_0\leq c_E^{F}\leq c_0+\beta,\quad
H_g^{F}\geq c_0-\beta\frac{a}{C_\star v_\star}gU'(g)
\geq c_0-\beta(1+v_+/w)=\eta>0.\tag{12}
\end{equation}
Use the old bound $gU'(g)\leq v(1+v/w)$ and $a\leq C(v)$.
Thus the old conservative root domain remains sufficient, as do the
rate bound $g_E\leq R/c_0$ and the bounded-work conditional clearance
interface. No O-R profile, reference mass, cutoff, source or physical-time
update is changed by (10).

\newpage
\section*{5. What survives, what does not, and the decision boundary}
\textbf{Conditional preservation.} At $v=v_+$ all immediate-work values
and queued/free trace rules coincide, so the previously constructed complete
regular cohort at that advice remains the same conditional solution if
W-F is selected. This is an algebraic anchor, not a preferred action.
Local continuation uses the same strict-event arguments, subject to their
existing scope. The O-R positive-budget proof uses the unchanged admission
rate and profile bounds; storage, conservation and local TST remain conditional on
admissible departures. No all-action hybrid theorem is supplied here.

\textbf{Not inherited.} Advice derivatives away from the anchor, optimal
actions, information regret, applied MP6 and SUMO usefulness must be qualified
for the selected work law. The candidate's coefficient, representative
insertion speed and maneuver-versus-recovery attribution are not calibrated
by (10). The benefit in (11) must not be promoted to a universal empirical
law or used to select parameters for an interior optimum.

\textbf{Two viable interpretations.} W-A retains (2) with an explicit
direct desired-speed/saturation influence, requiring its own evidence.
W-F uses (10), omitting that extra influence and retaining actual-state
interaction. Both preserve the compact point architecture and its known
limitations. W-F is provisionally recommended for causal clarity; neither
is selected by this note.

\textbf{J4-W1.} This is the smallest Type 2 G3 choice about the work law,
with dependent G4 sensitivity and MP6 consequences. It does not reopen
H1 homogeneity, T1 storage, S-E-P, R-E, O-R, the observation contract or
choke-line accounting. Source/unit/algebra transcription is Type 1 under
the chosen interpretation. A new independent speed state or spatial zone
would need a separate reopening; none is indicated by this finding.

\textbf{Evidence obtained.} Exact identities (4)--(5), static rational
diagnostic (6)--(8), monotone-root signs (9), candidate signs/bounds
(11)--(12), negative controls and anchor equality. Symbolic checks verify
identities and rational inequalities; they do not mechanically prove
physical validity, H1 reachability or every coupled event sequence.

\textbf{Next work after selection.} Resume J4 at whole original action-
range/event and physical-domain qualification; then the sufficient-information,
full-objective sensitivity, strong-rival and S-P--S-L contracts. Step 3 is
open; Step 4 and implementation/experiments are unopened.

\textbf{Primary references and reading limits.}
\begin{itemize}
\item \href{https://repository.tudelft.nl/file/File_c424c1b3-7374-436b-a6a7-abdf69063aa6}{Leclercq et al. (2016)}, Sections 2--3.2: restricted point/recovery context.
\item \href{https://mocopo.ifsttar.fr/fileadmin/contributeurs/Mocopo/Documents/Other_documents/Article_-_Merging_behaviour_-_Empirical_comparison_between_two_sites_and_new_theory_development.pdf}{Marczak et al. (2013)}, Sections 4.4--7: gaps, relative speeds and selection limits.
\item \href{https://arxiv.org/pdf/1405.7080}{Jin (2015)}, Sections 3 and 6: separate LTM boundary/node qualification.
\end{itemize}
No simulator, traffic time stepping, ODE rollout, numerical gradient,
parameter calibration or optimization is involved.
\end{document}
